% Patch 0624 A5_F^(4): Fourth-Order Closed-Form Substrate-Current Coefficients
% at Face-Aligned Reading C in the 600-Cell -- Empirical Validation of
% THEO-DSL-11 Theorem 1 at k=4
%
% OPEN-FP-F1-5 Sequence-4B closure artifact (THEO-DSL-12 candidate). Computes
% the O(delta^4) face-aligned coefficients alpha_4^(rho), alpha_4^(ax) in the
% 2D V_4-invariant subspace span{n_rho, n_ax} established at THEO-DSL-9 (Patch
% 0615; multi-AI confirmed at Patch 0617). EMPIRICALLY VALIDATES THEO-DSL-11
% Theorem 1 (the parity-dependent algebraic-structure theorem, Patches
% 0620/0622) at k=4 as a third independent data point.
%
% CHANGELOG
%   v1.0 (Patch 0624, Session 147, 28 May 2026): higher-order extension of
%     THEO-DSL-11 (face k=3, Patches 0620 + 0622 multi-AI confirmation).
%   v1.1 (Patch 0629, Session 147, 28 May 2026): consolidated wording-fix
%     revision pass per multi-AI review backlog from Patch 0626 reviews
%     (ChatGPT + Copilot) + Patch 0628 ChatGPT re-review CONFIRMED upgrade.
%     Specific changes:
%       - W1 (ChatGPT): replaced "genuine geometric law" framing with
%         "third independent empirical confirmation". The parity prediction
%         is *proven* at THEO-DSL-11; THEO-DSL-12 provides an empirical
%         stress-test data point.
%       - W2 (ChatGPT): added Remark (rem:thm-vs-validation) explicitly
%         distinguishing the proven THEO-DSL-11 Theorem 1 from the present
%         empirical stress-test artifact.
%       - W3 (ChatGPT): per-PSLQ-output wording adjusted from "exactly
%         zero" to "vanishing" sqrt(3) and sqrt(3)*phi slots. Exactness
%         belongs to the discovered integer relation, not to the
%         floating-point input.
%       - W5 (Copilot): added boxed Remark (rem:three-data-points)
%         explicitly listing the three-data-point empirical validation
%         of Theorem 1 across k=2,3,4 face variants.
%       - W6 (Copilot): clarified that Theorem 1 gives an upper bound on
%         the denominator (the ambient denominator), not an equality.
%       - W7 (Copilot): added cross-reference from main results section
%         to the vertex sign-alternation Observation block.
%     W4 (ChatGPT) is endorsement of current observational treatment for
%     sign-alternation; no change. No theorem changes; no claim weakened
%     or strengthened. All wording fixes preserve the artifact's
%     mathematical content; revisions are exposition-level only.
%
% ** SUBSTANTIVE NEW FINDING -- EMPIRICAL VALIDATION OF THEO-DSL-11 THEOREM 1 **:
%
% THEO-DSL-11 Theorem 1 predicts that face-aligned coefficients at order k
% live in:
%     - Q[phi]/3^(k/2)            at even k
%     - sqrt(3)*Q[phi]/3^((k+1)/2) at odd k
% with vertex and edge variants in Q[phi] at all k. At k=2 (even): clean
% Q[phi] confirmed by THEO-DSL-9 (-7/phi^2, 2 phi - 6). At k=3 (odd):
% sqrt(3)*Q[phi]/3 confirmed by THEO-DSL-11 ((87-53phi)sqrt(3)/3,
% (41-28phi)sqrt(3)/3).
%
% At k=4 (even), Theorem 1 predicts return to clean Q[phi]/9 ambient. The
% observed closed forms are TIGHTER than the ambient prediction:
%     alpha_4^(rho, face) = 641/2 - 180 phi      ~ +29.254  (denominator 2)
%     alpha_4^(ax)        = (401 - 167 phi)/3    ~ +43.596  (denominator 3)
% Both denominators 2 and 3 divide 9, consistent with the ambient Q[phi]/9
% prediction.
%
% PSLQ at mpmath dps=40 in extended basis {1, phi, sqrt(3), sqrt(3)*phi}
% returns integer relations [-2, 641, -360, 0, 0] and [-3, 401, -167, 0, 0]
% with vanishing sqrt(3) and sqrt(3)*phi slots. The k=4 face coefficients
% live in Q[phi], not in sqrt(3)*Q[phi] -- providing the third independent
% empirical confirmation of the parity-dependent algebraic structure
% predicted by THEO-DSL-11 Theorem 1, across k=2, 3, 4 face variants
% spanning both parities. (Theorem 1 is *proven* at THEO-DSL-11; this
% artifact provides an empirical stress-test data point, not a re-proof.)
%
% V_4-FORCED VANISHING of two non-invariant components persists at k=4
% (publication-grade Layer 3 unconditional from THEO-DSL-9 Lemmas 1-2, multi-
% AI confirmed at Patch 0617):
%     alpha_4^(perp3) = 0  (forced by sigma_F in V_4, all orders)
%     alpha_4^(diff)  = 0  (forced by sigma_E in V_4, all orders)
% Verified at machine zero across all 30 incident faces.
%
% Mandatory vertex cross-check at k=4: alpha_4^(vertex) = 855/2 - 252 phi
% ~ +19.755 (clean Q[phi]/2). Lower-order vertex cross-checks at k=1, 2, 3
% all reproduce (locking assembly convention).
%
% EMPIRICAL OBSERVATION (NOT A THEOREM): vertex sign-alternation conjecture
% empirically REFUTED at k=4. Sequence is +, -, +, + at k=1,2,3,4 (strict
% alternation broken). Registered as an observation-only programme finding,
% preserving honest scope-limitation.
%
% Path enumeration: 20,736 = 12^4 directed 4-edge paths. 30-face robustness
% std ~1e-13.
%
% Verify script: code/verify_face_alpha4_closure.py (ALL PASS checks).
% Reasoning fragment: hardened_theorems/reasoning/0624.md.

\documentclass[11pt]{article}
\usepackage[utf8]{inputenc}
\usepackage{amsmath,amssymb,amsthm}
\usepackage[margin=1in]{geometry}
\usepackage{hyperref}
\usepackage{booktabs}

\theoremstyle{plain}
\newtheorem{theorem}{Theorem}
\newtheorem{lemma}{Lemma}
\newtheorem{corollary}{Corollary}
\theoremstyle{definition}
\newtheorem*{remark}{Remark}
\newtheorem*{observation}{Empirical Observation}
\newtheorem*{antierasure}{Anti-erasure Remark}

\newcommand{\nhat}{\hat{n}}
\newcommand{\nrho}{\hat{n}_\rho}
\newcommand{\nax}{\hat{n}_{\text{ax}}}
\newcommand{\nperpF}{\hat{n}_{F\perp}}
\newcommand{\nperpthree}{\hat{n}_{\perp 3}}
\newcommand{\ndiff}{\hat{n}_{\text{diff}}}
\newcommand{\vhost}{v_{\text{host}}}
\newcommand{\ui}{u_i}
\newcommand{\uj}{u_j}
\newcommand{\jvec}{\vec{j}}

\title{Fourth-Order Closed-Form Substrate-Current Coefficients\\
at Face-Aligned Reading C in the 600-Cell\\
{\large \textbf{Empirical Validation of THEO-DSL-11 Theorem 1 at $k=4$}\\
(Sketch-Document Layer 3 under $(H5_F^{(4)})$; THEO-DSL-12 candidate;\\
OPEN-FP-F1-5 Sequence-4B; Patch 0624)}}
\author{CPP Programme}
\date{28 May 2026 (Session 147)}

\begin{document}
\maketitle

\begin{abstract}
This artifact closes Sequence-4B of OPEN-FP-F1-5, extending THEO-DSL-11 (the
face-aligned $\mathcal{O}(\delta^3)$ closure in $\sqrt{3}\cdot\mathbb{Q}[\phi]/3$,
multi-AI confirmed at Patches 0620 → 0622) to the next order in $\delta$.
Under the $V_4$-invariant 2-dimensional subspace $\mathrm{span}\{\nrho, \nax\}$
established at THEO-DSL-9, the order-four substrate-current vector at
face-aligned Reading C is given in closed form by
\begin{equation*}
\boxed{\;\;
  \jvec_4(\vhost) \;=\; \alpha_4^{(\rho)} \nrho + \alpha_4^{(\text{ax})} \nax,
  \quad
  \alpha_4^{(\rho)} = \frac{641}{2} - 180\phi,
  \quad
  \alpha_4^{(\text{ax})} = \frac{401 - 167\phi}{3}.
\;\;}
\end{equation*}
Numerically $\alpha_4^{(\rho)} \approx +29.254$ and $\alpha_4^{(\text{ax})}
\approx +43.596$. Both lie in clean $\mathbb{Q}[\phi]$ (denominators 2 and 3
respectively, both dividing the ambient $9 = 3^{k/2}$ predicted by
THEO-DSL-11 Theorem 1 at $k=4$).

\paragraph{Empirical validation of THEO-DSL-11 Theorem 1 at $k=4$ — the
central content of this artifact.}
The parity-dependent algebraic-structure theorem registered at THEO-DSL-11
predicts face-aligned coefficients live in $\mathbb{Q}[\phi]/3^{k/2}$ at
even $k$ and $\sqrt{3}\cdot\mathbb{Q}[\phi]/3^{(k+1)/2}$ at odd $k$. At
$k=4$ (even), this predicts return to clean $\mathbb{Q}[\phi]/9$ ambient.
PSLQ at mpmath dps=40 in the extended basis $\{1, \phi, \sqrt{3}, \sqrt{3}\phi\}$
returns the integer relations
\begin{align*}
[-2,\, 641,\, -360,\, 0,\, 0] &\quad\text{for}\quad \alpha_4^{(\rho)},\\
[-3,\, 401,\, -167,\, 0,\, 0] &\quad\text{for}\quad \alpha_4^{(\text{ax})}.
\end{align*}
\textbf{The $\sqrt{3}$ and $\sqrt{3}\phi$ slots of both integer relations
vanish identically.} This is the \textbf{third independent data point}
empirically validating Theorem 1, after THEO-DSL-9 at $k=2$ (clean $\mathbb{Q}[\phi]$)
and THEO-DSL-11 at $k=3$ (in $\sqrt{3}\cdot\mathbb{Q}[\phi]/3$). The
parity-dependent algebraic structure now has \textbf{three independent
empirical confirmations} spanning both parities ($k \in \{2, 3, 4\}$),
sharply distinguishing the parity prediction from a $k=3$ accident.

\paragraph{Structural inheritance preserved at $k=4$.}
The $V_4$-forced vanishing of the two non-invariant components carries
over to $k=4$ unconditionally from THEO-DSL-9 Lemmas 1--2 (multi-AI
confirmed at Patch 0617): $\alpha_4^{(\perp 3)} = 0$ (forced by $\sigma_F$)
and $\alpha_4^{(\text{diff})} = 0$ (forced by $\sigma_E$). Verified at
machine zero across all 30 incident faces, robustness std $\sim 10^{-13}$.

The vertex-aligned cross-check at $k=4$ reproduces $\alpha_4^{(\text{vertex})}
= 855/2 - 252\phi \approx +19.755$ (clean $\mathbb{Q}[\phi]/2$); lower-order
vertex cross-checks at $k=1, 2, 3$ all reproduce, locking the assembly
convention across $k \in \{1, 2, 3, 4\}$. The result is verified by
independent numerical assembly of all $20{,}736 = 12^4$ directed 4-edge
paths from the host vertex. The full computation is sketch-document
Layer 3 under the natural extension $(H5_F^{(4)})$ of the path-class weight
ansatz $W(P) = 1$ from $\mathcal{P}_3(\vhost)$ to $\mathcal{P}_4(\vhost)$.
\end{abstract}

\section{Scope and provenance}\label{sec:scope}

This artifact registers a new candidate hardened theorem THEO-DSL-12:
the fourth-order ($\mathcal{O}(\delta^4)$) closed-form substrate-current
coefficients at face-aligned Reading C in the 600-cell, under the
2-dimensional $V_4$-invariant subspace established at THEO-DSL-9 (Patch
0615; multi-AI confirmed at Patch 0617). It is the direct higher-order
continuation of THEO-DSL-11 (face $k=3$, Patches 0620 → 0622) and tests
the parity prediction of THEO-DSL-11 Theorem 1 (the parity-dependent
algebraic-structure theorem) at $k=4$.

\paragraph{Inheritance from prior theorems.}
\begin{itemize}
\item \textbf{THEO-DSL-9} (Patch 0615, multi-AI confirmed Patch 0617):
  face-aligned $\mathcal{O}(\delta^2)$ closure $\alpha_2^{(\rho)} = -7/\phi^2
  = -14 + 7\phi$ and $\alpha_2^{(\text{ax})} = -6 + 2\phi$, in clean
  $\mathbb{Q}[\phi]$; $V_4$/2D structural Lemmas 1--2 establishing the
  $V_4$-forced vanishing of $\hat n_{\perp 3}$ and $\hat n_{\text{diff}}$
  components at all orders.
\item \textbf{THEO-DSL-11} (Patch 0620, multi-AI confirmed Patch 0622):
  face-aligned $\mathcal{O}(\delta^3)$ closure in $\sqrt{3}\cdot\mathbb{Q}[\phi]/3$
  with the parity-dependent algebraic-structure Theorem 1 (publication-grade
  Layer 3 unconditional, symmetry-only argument). The path-enumeration
  machinery (1728 paths, full 9-shell classification) is extended here at
  $k=4$ to 20,736 paths.
\item \textbf{THEO-DSL-10} (Patch 0618, multi-AI confirmed Patch 0623):
  edge-aligned $\mathcal{O}(\delta^3)$ closure in clean $\mathbb{Q}[\phi]$,
  confirming the $D_5$-edge variant carries no $\sqrt{3}$ at any $k$ (a
  data point supporting the parity-dependent structure being specific to
  face-aligned).
\end{itemize}

\paragraph{Hypotheses retained.}
$(H1_F)$ Reading C with host $\vhost$ and face-centroid anchor
$\nperpF = (\vhost + \ui + \uj)/(\phi\sqrt{3})$ for a chosen triangular face
$\{\vhost, \ui, \uj\}$ at $\vhost$ ($\ui, \uj \in S_1$ with $\ui \cdot \uj
= \phi/2$);
$(H2)$ Mechanism A (anisotropic-rate path-integral first-order substrate-
current ansatz);
$(H3_F)$ residual symmetry $V_4 = \{e, \sigma_E, \sigma_F, C_2\}$ at $\vhost$,
multi-AI confirmed at Patch 0617;
$(H4)$ the 600-cell unit-circumradius edge graph;
$(H5_F^{(4)})$ the path-class weight $W(P) = 1$ extended from $\mathcal{P}_3(\vhost)$
to $\mathcal{P}_4(\vhost)$ (no $k$-dependent reweighting).

The closed-form coefficient result is sketch-document Layer 3 under
$(H5_F^{(4)})$. The $V_4$/2D structural sub-result is publication-grade
Layer 3 unconditional (inherited from THEO-DSL-9 Lemmas 1--2). The
parity-dependent algebraic-structure inheritance is publication-grade Layer 3
unconditional (inherited from THEO-DSL-11 Theorem 1, symmetry-only argument).

\section{The path-integral assembly at $k=4$}\label{sec:assembly}

The assembly is the natural extension of the $k=3$ assembly to four edges:
\begin{equation}\label{eq:j4-assembly}
\jvec_4(\vhost) \;=\; \sum_{P \in \mathcal{P}_4(\vhost)}
  \prod_{j=1}^{4} (\hat e_j \cdot \nperpF) \cdot \hat e_1,
\end{equation}
where the sum is over directed 4-edge paths $\vhost \to u_1 \to u_2 \to u_3
\to u_4$ and $\hat e_j$ is the unit edge direction of step $j$. The path
count is $|\mathcal{P}_4(\vhost)| = 12^4 = 20{,}736$, since the 600-cell is
12-regular.

\subsection{Parity prediction at $k=4$}

By THEO-DSL-11 Theorem 1, the face-aligned $k=4$ coefficients lie in
$\mathbb{Q}[\phi]/3^{4/2} = \mathbb{Q}[\phi]/9$ at the ambient level: each
factor $(\hat e_j \cdot \nperpF) = c_j/\sqrt{3}$ with $c_j \in \mathbb{Q}[\phi]$,
so the four-fold product
\begin{equation*}
\prod_{j=1}^{4} (\hat e_j \cdot \nperpF) \;=\; \frac{\prod_j c_j}{(\sqrt{3})^4}
   \;=\; \frac{\prod_j c_j}{9}
\end{equation*}
lives in $\mathbb{Q}[\phi]/9$. Summing over $12^4 = 20{,}736$ paths preserves
this, and projecting onto a basis vector in $\mathbb{Q}[\phi]^4$ (such as
$\nrho$ or $\nax$) yields a coefficient in $\mathbb{Q}[\phi]/9$.

This is an \textit{ambient upper bound} on the denominator: Theorem 1
bounds the denominator above by $3^{k/2} = 9$ at $k=4$ but does
\textbf{not} predict minimal denominators. The actual coefficient may
have a tighter denominator if particular sub-sums happen to cancel. As
we will see, both observed $k=4$ coefficients lie in $\mathbb{Q}[\phi]/9$
but with the tighter denominators 2 and 3 respectively (both dividing 9),
illustrating that Theorem 1's denominator is an upper bound rather than
an equality.

\subsection{Algebraic falsifiability}

The parity prediction is sharply falsifiable: if Theorem 1 were wrong (e.g.,
if the parity structure were an accident specific to $k=3$), one might
expect non-trivial $\sqrt{3}$ contributions at $k=4$. PSLQ in the extended
basis $\{1, \phi, \sqrt{3}, \sqrt{3}\phi\}$ at sufficiently high precision
would then return non-zero $\sqrt{3}$ and $\sqrt{3}\phi$ coefficients. The
present artifact's central empirical result is that PSLQ at mpmath dps=40
returns integer relations with \textbf{vanishing} $\sqrt{3}$ and
$\sqrt{3}\phi$ slots for both $\alpha_4^{(\rho)}$ and $\alpha_4^{(\text{ax})}$,
confirming the prediction at $k=4$ as a third independent empirical data
point for Theorem 1.

\section{Closed forms at $k=4$}\label{sec:results}

\begin{theorem}[THEO-DSL-12 candidate: face-aligned $\mathcal{O}(\delta^4)$
coefficients]\label{thm:dsl12}
Under hypotheses $(H1_F)$--$(H4)$ and $(H5_F^{(4)})$,
\begin{align}
\jvec_4(\vhost) &= \alpha_4^{(\rho)} \nrho + \alpha_4^{(\text{ax})} \nax, \\
\alpha_4^{(\rho)} &= \frac{641}{2} - 180\phi \approx +29.254, \\
\alpha_4^{(\text{ax})} &= \frac{401 - 167\phi}{3} \approx +43.596.
\end{align}
Both coefficients lie in $\mathbb{Q}[\phi]$ (denominators 2 and 3
respectively, both dividing the ambient $\mathbb{Q}[\phi]/9$ predicted by
THEO-DSL-11 Theorem 1). The two non-invariant components vanish identically:
\begin{equation}
\alpha_4^{(\perp 3)} = 0, \quad \alpha_4^{(\text{diff})} = 0,
\end{equation}
forced unconditionally by $V_4$-equivariance (THEO-DSL-9 Lemmas 1--2,
multi-AI confirmed at Patch 0617). The coefficients are independent of the
choice of incident face $\{\ui, \uj\}$ (verified explicitly over all 30
triangular faces incident to $\vhost$, robustness std $\sim 10^{-13}$).
\end{theorem}

\begin{proof}[Computation and verification]
The vector~\eqref{eq:j4-assembly} is evaluated by enumerating all
$20{,}736 = 12^4$ directed 4-edge paths from $\vhost$ in the 600-cell and
summing the product of four dot products against $\hat e_1$ (the first-edge
direction). The $V_4$-forced direction of the result --- $\jvec_4 \in
\mathrm{span}\{\nrho, \nax\}$ --- is the same structural lemma used at
$k = 2, 3$ (THEO-DSL-9 \S2.1, THEO-DSL-11 \S2) and carries over to $k=4$
unconditionally. The non-zero components in the basis $\{\nrho, \nax\}$
are extracted via inner products in the exact-unit-norm frame (Lemma 1 of
THEO-DSL-11 for $\nax$).

The closed forms in $\mathbb{Q}[\phi]/9$ ambient are identified via mpmath
PSLQ at dps=40 in the extended basis $\{1, \phi, \sqrt{3}, \sqrt{3}\phi\}$;
the relations
\begin{equation*}
[-2,\, 641,\, -360,\, 0,\, 0] \text{ for } \alpha_4^{(\rho)},
\qquad
[-3,\, 401,\, -167,\, 0,\, 0] \text{ for } \alpha_4^{(\text{ax})}
\end{equation*}
yield the closed forms $\alpha_4^{(\rho)} = (641 - 360\phi)/2 = 641/2 - 180\phi$
and $\alpha_4^{(\text{ax})} = (401 - 167\phi)/3$. The closed-form values are
then verified to machine-double precision in the verify script
\texttt{code/verify\_face\_alpha4\_closure.py} (all PASS checks).
\end{proof}

\begin{corollary}[empirical validation of THEO-DSL-11 Theorem 1 at $k=4$]
\label{cor:thm1-k4}
The PSLQ relations of Theorem~\ref{thm:dsl12} have \textbf{vanishing
$\sqrt{3}$ and $\sqrt{3}\phi$ slots}, confirming THEO-DSL-11 Theorem 1's
prediction that face-aligned coefficients live in $\mathbb{Q}[\phi]/3^{k/2}$
at even $k$. Combined with THEO-DSL-9 ($k=2$, clean $\mathbb{Q}[\phi]$) and
THEO-DSL-11 ($k=3$, $\sqrt{3}\cdot\mathbb{Q}[\phi]/3$), this is the third
independent empirical confirmation of Theorem 1's parity prediction,
spanning both parities in $k \in \{2, 3, 4\}$.
\end{corollary}

\begin{remark}[theorem status vs empirical validation status]\label{rem:thm-vs-validation}
THEO-DSL-11 Theorem 1 (the parity-dependent algebraic-structure theorem,
Patches 0620 + 0622) is \emph{proven} at publication-grade Layer 3
unconditional via symmetry + algebraic-home reasoning. The present
artifact (THEO-DSL-12 candidate) does \emph{not} re-prove Theorem 1;
it provides an independent empirical \emph{stress-test} of the
prediction at the next even $k$, complementing the existing $k=2$
(THEO-DSL-9) and $k=3$ (THEO-DSL-11) data points. The closed forms
$\alpha_4^{(\rho)} = 641/2 - 180\phi$ and $\alpha_4^{(\text{ax})} =
(401 - 167\phi)/3$ are themselves a new closed-form result; their
\emph{compatibility with Theorem 1's parity prediction} (i.e., the
$\sqrt{3}$ slots of the integer PSLQ relations vanishing) is the
stress-test outcome, not a re-proof of Theorem 1 itself.
\end{remark}

\begin{remark}[three-data-point empirical validation, summary]\label{rem:three-data-points}
\fbox{\parbox{0.92\textwidth}{
\textbf{Three-data-point validation of THEO-DSL-11 Theorem 1.}
The parity-dependent algebraic-structure theorem (THEO-DSL-11 Theorem 1)
is now empirically validated across $k \in \{2, 3, 4\}$ face-aligned
data points:
\begin{itemize}
\item $k=2$ (even, THEO-DSL-9): $\alpha_2^{(\rho)} = -14 + 7\phi$,
      $\alpha_2^{(\text{ax})} = -6 + 2\phi$, both in clean $\mathbb{Q}[\phi]$.
\item $k=3$ (odd, THEO-DSL-11): $\alpha_3^{(\rho)} = (87 - 53\phi)\sqrt{3}/3$,
      $\alpha_3^{(\text{ax})} = (41 - 28\phi)\sqrt{3}/3$, both in
      $\sqrt{3}\cdot\mathbb{Q}[\phi]/3$.
\item $k=4$ (even, this artifact): $\alpha_4^{(\rho)} = 641/2 - 180\phi$,
      $\alpha_4^{(\text{ax})} = (401 - 167\phi)/3$, both in $\mathbb{Q}[\phi]/9$
      ambient (with tighter denominators 2 and 3).
\end{itemize}
The parity-dependent algebraic-class membership $\mathbb{Q}[\phi]$ vs
$\sqrt{3}\cdot\mathbb{Q}[\phi]$ is confirmed empirically across all three data
points spanning both parities. No $\sqrt{3}$ slot has appeared at any even $k$
where Theorem 1 predicts clean $\mathbb{Q}[\phi]$, and the $\sqrt{3}$
prefactor at $k=3$ matches Theorem 1's odd-$k$ prediction exactly. The
appropriate epistemic framing of this artifact's central contribution is
\emph{third independent empirical confirmation of the parity prediction},
not establishment of a new theorem.
}}
\end{remark}

\begin{corollary}[$k=4$ vertex cross-check]\label{cor:vertex-k4}
The same 20,736-path assembly with $\nhat = \vhost$ (vertex-aligned cross-check)
reproduces $\alpha_4^{(\text{vertex})} = 855/2 - 252\phi \approx +19.755$, in
clean $\mathbb{Q}[\phi]/2$. The $k=1, 2, 3$ vertex cross-checks at $\alpha_1
= 3/\phi^2$, $\alpha_2 = -9/\phi^2$, $\alpha_3 = -126 + 81\phi$ are also all
reproduced, locking the assembly convention across $k \in \{1, 2, 3, 4\}$.
\end{corollary}

\noindent The vertex-aligned coefficients reproduced in
Corollary~\ref{cor:vertex-k4} additionally support a programme-finding
about the sign pattern across $k$, registered as an observation-only
result in Observation~\ref{obs:sign-alternation} below.

\begin{observation}[vertex sign-alternation conjecture refuted at $k=4$]
\label{obs:sign-alternation}
The sequence of vertex-aligned coefficients $(\alpha_1, \alpha_2, \alpha_3,
\alpha_4) \approx (+1.146, -3.438, +5.061, +19.755)$ has signs $(+, -, +, +)$
at $k = 1, 2, 3, 4$. \textbf{Strict alternation is empirically refuted at
$k=4$.} This observation refutes a programme conjecture that had been carried
since the THEO-DSL-3,-4,-5 vertex-aligned theorems and that ChatGPT had
endorsed keeping observational at Patch 0623. The refutation is registered
here as a programme finding; the underlying geometric reason for the
sign-pattern transition at $k=4$ is not yet structurally identified and is
flagged as a future investigative target. This is an \textbf{observation-only}
finding, NOT a theorem; no claim is made about higher-$k$ sign patterns
based on this single transition.
\end{observation}

\section{Robustness over face choice}\label{sec:robustness}

By the $V_4$ structural inheritance from THEO-DSL-9, the coefficient pair
$(\alpha_4^{(\rho)}, \alpha_4^{(\text{ax})})$ is invariant under the
$H_4$-orbit action on the 30 incident faces $\{\ui, \uj\}$ at $\vhost$, when
expressed in the corresponding face-adapted $\{\nrho, \nax\}$ frame (per the
ChatGPT-0622 clarification: the frame rotates with face choice, but the
projected coefficients in the rotated frame are orbit-invariant). The verify
script confirms this explicitly: across all 30 faces, the standard deviation
of each coefficient over the orbit is $\sim 10^{-13}$ (double precision);
the $\hat n_{\perp 3}$ and $\hat n_{\text{diff}}$ components are identically
zero to $\sim 10^{-13}$.

\section{Exclusions retained}\label{sec:exclusions}
\begin{enumerate}
\item[$(E1)$] No path-class weight other than $W(P) = 1$ at $k=4$ is
  considered; $(H5_F^{(4)})$ is the natural extension of $(H5_F^{(3)})$
  used at THEO-DSL-11.
\item[$(E2)$] Vertex (THEO-DSL-3,-4,-5; this artifact's cross-check) and edge
  (THEO-DSL-7,-10) variants are separate artifacts at the same order;
  THEO-DSL-12 is the face-aligned $k=4$ artifact only. The edge-aligned
  $k=4$ closure (Sequence-4A) is a future OPEN-FP-F1-5 target.
\item[$(E3)$] No promotion of Layer 3 status to publication-grade for the
  coefficient closure: the underlying ansatz $(H5_F^{(4)})$ is sketch-document
  Layer 3. The $V_4$/2D structural sub-results (vanishing of $\alpha_4^{(\perp 3)}$
  and $\alpha_4^{(\text{diff})}$) and the parity-dependent algebraic-structure
  inheritance from THEO-DSL-11 Theorem 1 are publication-grade Layer 3
  unconditional.
\item[$(E4)$] No claim that the parity-dependent $\mathbb{Q}[\phi]$ vs
  $\sqrt{3}\cdot\mathbb{Q}[\phi]$ structure (THEO-DSL-11 Theorem 1) extends
  to other Reading-C variants beyond face-aligned. The vertex and edge
  variants stay in $\mathbb{Q}[\phi]$ at all $k$ since their substrate
  primitive directions ($\nrho = \vhost$, $\hat n_{\text{edge}}$) carry no
  $\sqrt{3}$.
\item[$(E5)$] No extension to $k \geq 5$ in this artifact. Theorem 1
  predicts $k=5$ face-aligned in $\sqrt{3}\cdot\mathbb{Q}[\phi]/9$ (odd $k$);
  $k=6$ face-aligned in $\mathbb{Q}[\phi]/27$ (even $k$); etc. These
  predictions are not yet computationally verified.
\item[$(E6)$] The vertex sign-alternation observation (Observation~\ref{obs:sign-alternation})
  is empirical only; no structural theorem is offered for the sign pattern
  at $k \geq 5$.
\end{enumerate}

\section{Provenance and verification}\label{sec:provenance}

The exact closed forms above are obtained by:
\begin{enumerate}
\item Building the 600-cell from first principles (same construction as
  THEO-DSL-9/-10/-11 verify scripts) and verifying its 120-vertex,
  12-regular structure.
\item Verifying the exact-norm identities $\|\nax\| = 1$ and $\|\vhost + \ui
  + \uj\| = \phi\sqrt{3}$ (Lemmas 1--2 of THEO-DSL-11, reproduced here).
\item Enumerating all $20{,}736 = 12^4$ directed 4-edge paths $\vhost \to
  u_1 \to u_2 \to u_3 \to u_4$ and computing the assembly~\eqref{eq:j4-assembly}
  with full 9-shell vertex classification (the 4-shell coarsening that
  suffices at $k \leq 2$ does NOT suffice at $k \geq 3$; cf. THEO-DSL-10
  reasoning fragment 0618.md \S6).
\item Decomposing the result in the orthonormal basis $\{\nrho, \nax,
  \nperpthree, \ndiff\}$ via direct dot products.
\item Extracting closed forms via mpmath PSLQ at dps=40 in the extended
  basis $\{1, \phi, \sqrt{3}, \sqrt{3}\phi\}$; the integer relations
  $[-2, 641, -360, 0, 0]$ and $[-3, 401, -167, 0, 0]$ confirm the closed
  forms with vanishing $\sqrt{3}$ and $\sqrt{3}\phi$ slots.
\item Verifying the resulting closed-form expressions to machine-double
  precision in the verify script.
\item Cross-checking via the vertex-aligned $k=1, 2, 3, 4$ targets
  ($3/\phi^2$, $-9/\phi^2$, $-126 + 81\phi$, $855/2 - 252\phi$), the
  THEO-DSL-9 reproduction at $k=2$ face, the THEO-DSL-11 reproduction at
  $k=3$ face, and the $V_4$-forced vanishing of the two non-invariant
  components (machine zero across all 30 faces).
\end{enumerate}
See \texttt{code/verify\_face\_alpha4\_closure.py} for the full verification.
The verbatim reasoning trail is captured in \texttt{hardened\_theorems/reasoning/0624.md}.

\begin{antierasure}
The empirical validation of THEO-DSL-11 Theorem 1 at $k=4$ across the
parity-dependent structure $(k=2 \text{ clean}, k=3 \text{ with } \sqrt{3},
k=4 \text{ clean again})$ is the central content of this artifact. Should
a future multi-AI review surface either (i) a counter-derivation yielding
non-zero $\sqrt{3}$ contributions at $k=4$ (which would refute Theorem 1's
parity prediction), (ii) coefficient values different from those stated
in Theorem~\ref{thm:dsl12}, (iii) a $V_4$-forced-vanishing failure at $k=4$,
or (iv) a vertex cross-check failure at $k=4$, this artifact must be
retained verbatim and a separate revision/supersession artifact created.
The vertex sign-alternation conjecture refutation (Observation~\ref{obs:sign-alternation})
is registered as observation-only; no claim is made about $k \geq 5$ sign
patterns. Multi-AI verification is invited on the Theorem 1 empirical
validation at $k=4$, paralleling the THEO-DSL-9 review cycle (Patches
0615 → 0617) and the THEO-DSL-11 review cycle (Patches 0620 → 0622).
\end{antierasure}

\end{document}
